% ============================================================
% Informe L3 - Redes de Computadores - Grupo 14
% Compilar con pdfLaTeX (Overleaf: Compiler = pdfLaTeX, Main document = main.tex)
% Cada seccion vive en su propio archivo para que los tres puedan
% editar a la vez sin pisarse (ver README.md).
% ============================================================
\documentclass[11pt,a4paper]{article}

\usepackage[utf8]{inputenc}
\usepackage[T1]{fontenc}
\usepackage[spanish,es-tabla]{babel}
\usepackage{lmodern}
\usepackage[margin=2.3cm]{geometry}
\usepackage{graphicx}
\usepackage{booktabs,longtable,array,tabularx}
\usepackage{xcolor}
\usepackage{amsmath}
\usepackage{listings}
\usepackage{enumitem}
\usepackage{caption}
\usepackage{fancyhdr}
\usepackage[hidelinks]{hyperref}

\graphicspath{{evidencias/}}
\setlength{\parindent}{0pt}
\setlength{\parskip}{0.6em}
\captionsetup{font=small,labelfont=bf}

% ---------- encabezado ----------
\setlength{\headheight}{14pt}
\pagestyle{fancy}
\fancyhf{}
\lhead{Redes de Computadores -- Laboratorio 3}
\rhead{Grupo 14}
\cfoot{\thepage}

% ---------- variables del contrato (generadas) ----------
\newcommand{\Xval}{4}
\newcommand{\Yval}{14}
\newcommand{\CriterioSubredes}{2}


% ---------- macros utiles ----------
% Texto pendiente: sale en rojo, para no entregar nada sin revisar.
% Antes de entregar: buscar "PENDIENTE" en el PDF y en el codigo.
\newcommand{\pend}[1]{\textcolor{red}{[PENDIENTE: #1]}}
% Marca de autor al inicio de cada parte.
\newcommand{\autor}[1]{\textit{Responsable: #1}\par}
% Comando de consola en linea.
\newcommand{\cmd}[1]{\texttt{#1}}
% Captura de Packet Tracer: \evid{ruta/relativa/a/evidencias.png}{leyenda}
% Si el archivo aun no existe, deja un recuadro visible en lugar de fallar.
\newcommand{\evid}[2]{%
  \IfFileExists{evidencias/#1}{%
    \begin{figure}[htbp]\centering
      \includegraphics[width=0.92\linewidth]{#1}
      \caption{#2}
    \end{figure}}{%
    \begin{center}\fbox{\parbox{0.9\linewidth}{\centering
      \textcolor{red}{Falta captura \texttt{\detokenize{evidencias/#1}}}\\ #2}}\end{center}}}

% ---------- listados de configuracion IOS ----------
\lstdefinestyle{ios}{
  basicstyle=\ttfamily\scriptsize,
  commentstyle=\color{gray},
  breaklines=true,
  frame=single,
  rulecolor=\color{black!30},
  columns=fullflexible,
  keepspaces=true,
  showstringspaces=false,
  inputencoding=utf8,
  extendedchars=true,
  morecomment=[l]{!},
  literate={á}{{\'a}}1 {é}{{\'e}}1 {í}{{\'i}}1 {ó}{{\'o}}1 {ú}{{\'u}}1
           {Á}{{\'A}}1 {É}{{\'E}}1 {Í}{{\'I}}1 {Ó}{{\'O}}1 {Ú}{{\'U}}1
           {ñ}{{\~n}}1 {Ñ}{{\~N}}1,
}
\lstset{style=ios}

\begin{document}

\begin{titlepage}
  \centering
  {\large Universidad Técnica Federico Santa María\par}
  {\large Departamento de Informática\par}
  \vspace{1.5cm}
  {\Large Redes de Computadores\par}
  \vspace{0.8cm}
  {\huge\bfseries Laboratorio 3\par}
  {\Large Capa de Red: simulación en Cisco Packet Tracer\par}
  \vspace{1.2cm}
  {\Large Grupo 14\par}
  \vspace{1.5cm}

  \begin{tabular}{@{}ll@{}}
    \toprule
    \textbf{Integrante} & \textbf{Rol USM} \\
    \midrule
    Sebastián \pend{apellidos} & \pend{rol USM} \\
    Benjamín \pend{apellidos}  & \pend{rol USM} \\
    Bárbara \pend{apellidos}   & \pend{rol USM} \\
    \bottomrule
  \end{tabular}

  \vspace{1.5cm}
  {\large Profesor: Jorge Díaz\par}
  {\large Ayudantes: Cristóbal Brandau, Juan Cucurella\par}
  \vfill
  {\large Octubre de 2026\par}
\end{titlepage}


\tableofcontents
\newpage

% La portada es la seccion 1 del enunciado; el resto parte en 2.
\setcounter{section}{1}

\section{Memoria de cálculo de subredes}

\autor{cada integrante completa sus LAN; las tablas se generan con \cmd{scripts/generar.py}}

\subsection{Sustitución de variables}

El número de grupo es $Y=\Yval$ y su último dígito es $X=\Xval$. Con ellos, las reglas del enunciado quedan así:

\begin{center}
\begin{tabular}{@{}p{2cm}p{9cm}p{1.5cm}@{}}
\toprule
LAN & Regla aplicada (X=\Xval, Y=\Yval) & Prefijo \\
\midrule
Deadlock & 172.16.0.0/16 en 32 subredes, la 3.ª & /21 \\
Raze & 172.16.0.0/16 en 32 subredes, la 8.ª & /21 \\
Killjoy & 172.16.0.0/16 en 32 subredes, la 15.ª & /21 \\
Cypher & 172.16.0.0/16 en 32 subredes, la 24.ª & /21 \\
Gekko & 20.25.0.0/16 con (X/2)+2 = 4 bits, la primera & /20 \\
Viper & 48.7.0.0/16 con Y+25 = 39 hosts, la primera & /26 \\
Reyna & 183.0.0.0/8 en al menos 2(X+8) = 24 subredes, la Y-ésima & /13 \\
Omen & 4.93.0.0/16 en 32 subredes, la Y-ésima & /21 \\
Brimstone & 192.168.50.0/24 (dada por el enunciado) & /24 \\
\bottomrule
\end{tabular}

\end{center}

Criterio de conteo de subredes adoptado: \textbf{criterio \CriterioSubredes}
(\pend{explicar el criterio elegido y por qué: 1 = la subred cero cuenta como primera; 2 = se excluye la subred cero}).

\subsection{Subredes LAN}

{\footnotesize\setlength{\tabcolsep}{3pt}
\begin{longtable}{@{}>{\raggedright\arraybackslash}p{1.6cm}>{\raggedright\arraybackslash}p{2.4cm}>{\raggedright\arraybackslash}p{2.5cm}>{\raggedright\arraybackslash}p{3.1cm}>{\raggedright\arraybackslash}p{2.2cm}>{\raggedright\arraybackslash}p{3.5cm}@{}}
\toprule
Subred & Red & Máscara y CIDR & Rango útil & Broadcast & Dispositivos \\
\midrule
\endhead
LAN Deadlock & 172.16.24.0/21 & 255.255.248.0 /21 & 172.16.24.1 -- 172.16.31.254 & 172.16.31.255 & Router Deadlock (172.16.24.1); GravNet (172.16.24.2), Annihilation (172.16.24.3) \\
LAN Raze & 172.16.64.0/21 & 255.255.248.0 /21 & 172.16.64.1 -- 172.16.71.254 & 172.16.71.255 & Router Raze (172.16.64.1); Paint Shells (172.16.64.2), Pack Blast (172.16.64.3) \\
LAN Killjoy & 172.16.120.0/21 & 255.255.248.0 /21 & 172.16.120.1 -- 172.16.127.254 & 172.16.127.255 & Router Killjoy (172.16.120.1); Turret (172.16.120.2), Lockdown (172.16.120.3) \\
LAN Cypher & 172.16.192.0/21 & 255.255.248.0 /21 & 172.16.192.1 -- 172.16.199.254 & 172.16.199.255 & Router Cypher (172.16.192.1); Trapwire (172.16.192.2), Spycam (172.16.192.3), Base server (172.16.192.4) \\
LAN Gekko & 20.25.16.0/20 & 255.255.240.0 /20 & 20.25.16.1 -- 20.25.31.254 & 20.25.31.255 & Router Gekko (20.25.16.1); Wingman, Dizzy, Thrash (DHCP) \\
LAN Viper & 48.7.0.64/26 & 255.255.255.192 /26 & 48.7.0.65 -- 48.7.0.126 & 48.7.0.127 & Router Viper (48.7.0.65); Poison Cloud, Snake Bite, Viper's Pit (DHCP) \\
LAN Reyna & 183.112.0.0/13 & 255.248.0.0 /13 & 183.112.0.1 -- 183.119.255.254 & 183.119.255.255 & Router Reyna (183.112.0.1); Dismiss (183.112.0.2), Devour (183.112.0.3), Empress (183.112.0.4) \\
LAN Omen & 4.93.112.0/21 & 255.255.248.0 /21 & 4.93.112.1 -- 4.93.119.254 & 4.93.119.255 & Router Omen (4.93.112.1); Paranoia (4.93.112.2), Dark cover (4.93.112.3), Shrouded Step (4.93.112.4) \\
LAN Brimstone & 192.168.50.0/24 & 255.255.255.0 /24 & 192.168.50.1 -- 192.168.50.254 & 192.168.50.255 & Router Brimstone (192.168.50.1); Location server (192.168.50.2) \\
\bottomrule
\end{longtable}

}

\subsection{Enlaces entre routers}

El router nombrado primero en cada enlace toma la dirección \texttt{.1} y el segundo la \texttt{.2}.

{\footnotesize\setlength{\tabcolsep}{3pt}
\begin{longtable}{@{}>{\raggedright\arraybackslash}p{1.8cm}>{\raggedright\arraybackslash}p{2.4cm}>{\raggedright\arraybackslash}p{4.2cm}>{\raggedright\arraybackslash}p{2.6cm}>{\raggedright\arraybackslash}p{4.2cm}@{}}
\toprule
Red & Máscara y CIDR & Rango útil & Broadcast & Dispositivos \\
\midrule
\endhead
1.0.0.0/8 & 255.0.0.0 /8 & 1.0.0.1 -- 1.255.255.254 & 1.255.255.255 & Deadlock (1.0.0.1), Cypher (1.0.0.2) \\
2.0.0.0/8 & 255.0.0.0 /8 & 2.0.0.1 -- 2.255.255.254 & 2.255.255.255 & Deadlock (2.0.0.1), Raze (2.0.0.2) \\
3.0.0.0/8 & 255.0.0.0 /8 & 3.0.0.1 -- 3.255.255.254 & 3.255.255.255 & Cypher (3.0.0.1), Killjoy (3.0.0.2) \\
4.0.0.0/8 & 255.0.0.0 /8 & 4.0.0.1 -- 4.255.255.254 & 4.255.255.255 & Raze (4.0.0.1), Killjoy (4.0.0.2) \\
5.0.0.0/8 & 255.0.0.0 /8 & 5.0.0.1 -- 5.255.255.254 & 5.255.255.255 & Raze (5.0.0.1), Neon (5.0.0.2) \\
6.0.0.0/8 & 255.0.0.0 /8 & 6.0.0.1 -- 6.255.255.254 & 6.255.255.255 & Killjoy (6.0.0.1), Fade (6.0.0.2) \\
7.0.0.0/8 & 255.0.0.0 /8 & 7.0.0.1 -- 7.255.255.254 & 7.255.255.255 & Fade (7.0.0.1), Viper (7.0.0.2) \\
8.0.0.0/8 & 255.0.0.0 /8 & 8.0.0.1 -- 8.255.255.254 & 8.255.255.255 & Neon (8.0.0.1), Gekko (8.0.0.2) \\
9.0.0.0/8 & 255.0.0.0 /8 & 9.0.0.1 -- 9.255.255.254 & 9.255.255.255 & Viper (9.0.0.1), Omen (9.0.0.2) \\
10.0.0.0/8 & 255.0.0.0 /8 & 10.0.0.1 -- 10.255.255.254 & 10.255.255.255 & Omen (10.0.0.1), Sage (10.0.0.2) \\
11.0.0.0/8 & 255.0.0.0 /8 & 11.0.0.1 -- 11.255.255.254 & 11.255.255.255 & Reyna (11.0.0.1), Sage (11.0.0.2) \\
12.0.0.0/8 & 255.0.0.0 /8 & 12.0.0.1 -- 12.255.255.254 & 12.255.255.255 & Gekko (12.0.0.1), Reyna (12.0.0.2) \\
13.0.0.0/8 & 255.0.0.0 /8 & 13.0.0.1 -- 13.255.255.254 & 13.255.255.255 & Sage (13.0.0.1), Brimstone (13.0.0.2) \\
\bottomrule
\end{longtable}

}

\section{Registro de configuración}

Los bloques siguientes son los comandos aplicados en cada router, tal como están en la carpeta \texttt{configs/} del repositorio.
Se ordenan por sistema autónomo.

\subsection*{Supuestos de implementación}

La topología entregada contiene los dispositivos con líneas dibujadas y sin cables reales, y cada router 1941 trae solo dos puertos Gigabit. El grupo la tomó como una referencia y supuso lo siguiente:
\begin{itemize}[nosep]
  \item Los enlaces se cablearon según la tabla siguiente. Los enlaces seriales usan un módulo HWIC-2T en la ranura 0 de cada router que lo necesita, con \cmd{clock rate 64000} en el extremo DCE (el router nombrado primero).
  \item Cada LAN se conecta a la interfaz \texttt{GigabitEthernet0/0} de su router mediante el switch; Base server se conecta al switch de Cypher para quedar en la misma LAN que Tripwire y Spycam.
  \item Los nombres de los equipos son los que aparecen en la topología entregada.
  \item Packet Tracer no soporta iBGP (al configurar un vecino del mismo AS responde ``Only external neighbors are supported''). Por eso, dentro del AS 40, Reyna, Omen y Sage corren OSPF (proceso 40) entre ellos y cada uno redistribuye entre OSPF y BGP con \cmd{redistribute bgp 40 subnets} y \cmd{redistribute ospf 40 match internal external 1 external 2}.
  \item En Raze y Killjoy la redistribución de OSPF hacia BGP se hace con \cmd{redistribute ospf 1 match internal}. Sin el filtro, Packet Tracer reexporta hacia el AS 20 las rutas externas aprendidas de Omega y se forma un bucle entre Neon y Raze (ver pregunta 3).
\end{itemize}

\begin{center}\footnotesize
\begin{tabular}{@{}llll@{}}
\toprule
Red & Cable & Extremo .1 & Extremo .2 \\
\midrule
1.0.0.0/8 & Cobre cruzado & Deadlock GigabitEthernet0/1 & Cypher GigabitEthernet0/1 \\
2.0.0.0/8 & Serial DCE & Deadlock Serial0/0/0 & Raze Serial0/0/0 \\
3.0.0.0/8 & Serial DCE & Cypher Serial0/0/0 & Killjoy Serial0/0/0 \\
4.0.0.0/8 & Serial DCE & Raze Serial0/0/1 & Killjoy Serial0/0/1 \\
5.0.0.0/8 & Cobre cruzado & Raze GigabitEthernet0/1 & Neon GigabitEthernet0/0 \\
6.0.0.0/8 & Cobre cruzado & Killjoy GigabitEthernet0/1 & Fade GigabitEthernet0/0 \\
7.0.0.0/8 & Cobre cruzado & Fade GigabitEthernet0/1 & Viper GigabitEthernet0/1 \\
8.0.0.0/8 & Cobre cruzado & Neon GigabitEthernet0/1 & Gekko GigabitEthernet0/1 \\
9.0.0.0/8 & Serial DCE & Viper Serial0/0/0 & Omen Serial0/0/0 \\
10.0.0.0/8 & Cobre cruzado & Omen GigabitEthernet0/1 & Sage GigabitEthernet0/0 \\
11.0.0.0/8 & Cobre cruzado & Reyna GigabitEthernet0/1 & Sage GigabitEthernet0/1 \\
12.0.0.0/8 & Serial DCE & Gekko Serial0/0/0 & Reyna Serial0/0/0 \\
13.0.0.0/8 & Serial DCE & Sage Serial0/0/0 & Brimstone Serial0/0/0 \\
\bottomrule
\end{tabular}

\end{center}

\newcommand{\registro}[2]{%
  \subsection{#1 (AS #2)}
  \lstinputlisting[caption={Configuración de #1}]{configs/#1.txt}}

\subsection*{AS 10 -- Base Alfa (OSPF área 0)}
\registro{Deadlock}{10}
\registro{Cypher}{10}
\registro{Raze}{10}
\registro{Killjoy}{10}

\subsection*{AS 20 -- Frontera Omega}
\registro{Neon}{20}
\registro{Fade}{20}

\subsection*{AS 30 -- Soporte táctico (con DHCP)}
\registro{Gekko}{30}
\registro{Viper}{30}

\subsection*{AS 40 -- Defensa y control}
\registro{Reyna}{40}
\registro{Omen}{40}
\registro{Sage}{40}

\subsection*{AS 50 -- Comando supremo}
\registro{Brimstone}{50}


\section{Desarrollo de preguntas y análisis}
\subsection{Pregunta 1: análisis del proceso DHCP (DORA)}
\autor{Benjamín}

Se capturó en modo Simulación el intercambio entre el host Wingman y el router Gekko, que actúa como servidor DHCP de la LAN 20.25.16.0/20 (dirección del router: 20.25.16.1). Para forzar una solicitud nueva, el host se pasó a configuración estática y luego otra vez a DHCP.

\evid{Q1_DHCP/dora_simulacion.png}{Mensajes DORA entre Wingman y Gekko en modo Simulación.}

\subsubsection*{Los cuatro mensajes}

Los mensajes viajan en paquetes UDP: el cliente usa el puerto 68 y el servidor el 67. Wingman aún no tiene dirección IP durante todo el intercambio, por lo que todos van con destino de difusión (\texttt{255.255.255.255}) y el switch Gekko~SW los replica a todos sus puertos; Dizzy y Thrash los reciben, pero los descartan porque no son servidores ni los destinatarios de la configuración.

\begin{center}\footnotesize
\setlength{\tabcolsep}{4pt}
\begin{tabular}{@{}>{\raggedright\arraybackslash}p{2.3cm}>{\raggedright\arraybackslash}p{2.3cm}>{\raggedright\arraybackslash}p{1.9cm}>{\raggedright\arraybackslash}p{2.3cm}>{\raggedright\arraybackslash}p{1.6cm}>{\raggedright\arraybackslash}p{3.8cm}@{}}
  \toprule
  Mensaje & Sentido & IP origen & IP destino & Puertos UDP & Función \\
  \midrule
  Discover & Wingman $\rightarrow$ difusión & 0.0.0.0 & 255.255.255.255 & 68 $\rightarrow$ 67 & El cliente busca servidores DHCP en su segmento. \\
  Offer & Gekko $\rightarrow$ Wingman & 20.25.16.1 & 255.255.255.255 & 67 $\rightarrow$ 68 & El servidor propone una configuración de red. \\
  Request & Wingman $\rightarrow$ Gekko & 0.0.0.0 & 255.255.255.255 & 68 $\rightarrow$ 67 & El cliente acepta la oferta; en la capa de aplicación indica el servidor elegido (20.25.16.1). \\
  Acknowledgment & Gekko $\rightarrow$ Wingman & 20.25.16.1 & 255.255.255.255 & 67 $\rightarrow$ 68 & El servidor confirma la asignación. \\
  \bottomrule
\end{tabular}
\end{center}

Como resultado del intercambio, Wingman recibió la dirección \texttt{20.25.16.11}, la máscara \texttt{255.255.240.0} y la puerta de enlace \texttt{20.25.16.1}. La primera dirección asignada es la .11 porque las diez primeras direcciones útiles de la subred (de .1 a .10) están excluidas del pool para asignaciones estáticas.

\pend{confirmar el tipo de cada mensaje en la pestaña Inbound/Outbound PDU Details de Packet Tracer}

\evid{Q1_DHCP/offer_wingman.png}{Offer recibida por Wingman (capas 2, 3, 4 y 7).}
\evid{Q1_DHCP/request_wingman.png}{Wingman recibe la Offer y responde con el Request.}
\evid{Q1_DHCP/request_ack_gekko.png}{Gekko recibe el Request y responde con el Acknowledgment.}

En la simulación, la Offer aparece alrededor de 2 segundos (tiempo simulado) después del Discover.

\subsubsection*{Direcciones IP del Discover}

\evid{Q1_DHCP/discover_capa3.png}{Detalle del DHCP Discover en Gekko: capas 2, 3, 4 y 7.}

En la capa de red, el Discover lleva como \textbf{origen 0.0.0.0} y como \textbf{destino 255.255.255.255}. En la capa de enlace, la MAC de origen es la de Wingman (\texttt{00D0.BC6D.B529}) y la de destino es la de difusión (\texttt{FFFF.FFFF.FFFF}).

Se usan esas direcciones porque el cliente todavía no tiene ninguna dirección IP asignada, así que no puede usar una como origen y usa la dirección ``no especificada'' \texttt{0.0.0.0}. Tampoco conoce la dirección del servidor, por lo que envía el mensaje a la dirección de difusión limitada \texttt{255.255.255.255}, que llega a todos los equipos del segmento sin que un router la reenvíe a otras redes. Así, cualquier servidor DHCP presente en la LAN puede recibirlo y responder.

\subsubsection*{Segmento sin pool DHCP}

Para comprobarlo se eliminó temporalmente el pool en Gekko (\cmd{no ip dhcp pool POOL\_LAN}) y se repitió la solicitud desde Wingman. Packet Tracer informó que la solicitud DHCP falló, y el host quedó con la dirección \textbf{169.254.181.41}. Después se restauró el pool y Wingman volvió a obtener una dirección de la subred 20.25.16.0/20.

Al no haber ningún pool, el router no genera una Offer, por lo que el cliente nunca recibe configuración de red. Tras agotar sus intentos, el host se asigna a sí mismo una dirección del rango \texttt{169.254.0.0/16}, una dirección \textit{link-local} (APIPA) con máscara \texttt{255.255.0.0} \pend{verificar la máscara en la captura}. Este estado significa que la autoconfiguración falló: la dirección solo permite comunicarse con otros equipos del mismo segmento que también se hayan autoasignado una, y no incluye puerta de enlace, por lo que el host no puede salir de su LAN.

\evid{Q1_DHCP/sin_pool.png}{Host solicitando DHCP en un segmento sin pool.}

\subsection{Pregunta 2: verificación global de conectividad}
\autor{Bárbara}

\textbf{Pruebas pedidas:} ping y traceroute de GravNet a Location server, y de Spycam a Shrouded Step, con cada router intermedio indicado.

\subsubsection*{GravNet $\rightarrow$ Location server}
\evid{Q2_conectividad/ping_gravnet_location.png}{Ping de GravNet a Location server.}
\evid{Q2_conectividad/tracert_gravnet_location.png}{Traza de GravNet a Location server.}

\begin{tabular}{@{}cll@{}}
  \toprule
  Salto & Router & IP observada \\
  \midrule
  1 & Deadlock & 172.16.24.1 \\
  2 & Cypher & 1.0.0.2 \\
  3 & Neon & 5.0.0.2 \\
  4 & Gekko & 8.0.0.2 \\
  5 & Reyna & 12.0.0.2 \\
  6 & Sage & 11.0.0.2 \\
  7 & Brimstone & 13.0.0.2 \\
  8 & Location server & 192.168.50.2 \\
  \bottomrule
\end{tabular}

\medskip
El paquete recorre 7 routers y atraviesa los cinco sistemas autónomos de la red (AS 10, 20, 30, 40 y 50) hasta llegar al Location server. El ping responde
sin pérdidas y el TTL alterna entre 120 y 121, lo que indica que los paquetes no siguen siempre el mismo camino: algunos atraviesan un router más que otros.
Esto es coherente con el reparto del tráfico entre dos rutas de igual métrica que Packet Tracer muestra en Deadlock y Cypher (ver pregunta 3). Entre Cypher y Neon la traza no muestra
routers intermedios, porque en cada salto Packet Tracer informa la IP de la interfaz por la que responde y no por la que entra el paquete (Ver preguntas 3 y 4)


\subsubsection*{Spycam $\rightarrow$ Shrouded Step}
\evid{Q2_conectividad/ping_spycam_shrouded.png}{Ping de Spycam a Shrouded Step.}
\evid{Q2_conectividad/tracert_spycam_shrouded.png}{Traza de Spycam a Shrouded Step.}

\begin{tabular}{@{}cll@{}}
  \toprule
  Salto & Router & IP observada \\
  \midrule
  1 & Cypher & 172.16.192.1 \\
  2 & Killjoy & 3.0.0.2 \\
  3 & Fade & 6.0.0.2 \\
  4 & Viper & 7.0.0.2 \\
  5 & Omen & 9.0.0.2 \\
  6 & Shrouded Step & 4.93.112.4 \\
  \bottomrule
\end{tabular}

\medskip
El paquete sale del AS 10 por Killjoy, cruza Fade (AS 20) y Viper (AS 30) y entra al AS 40 por Omen, para llegar hasta Shrouded Step en 6 saltos. Killjoy prefiere esta salida porque 
aprende la red de la LAN de Omen por eBGP desde Fade, con distancia administrativa 20, mientras que la misma red también le llega por OSPF con distancia 110. El router instala la de 
menor distancia (Ver la tabla de Killjoy en la pregunta 6).
\subsection{Pregunta 3: comparativa de selección de rutas (OSPF vs.\ BGP)}
\autor{Sebastián}

Desde Turret (172.16.120.2, LAN de Killjoy) se envió un ping y una traza hacia Location server (192.168.50.2, AS 50). El ping respondió los cuatro paquetes con TTL 121, o sea que la respuesta pasó por siete routers antes de llegar a Turret.

\evid{Q3_ospf_vs_bgp/ping_tracert_turret_location.png}{Ping y traza de Turret a Location server.}

\subsubsection*{Ruta seguida}

El paquete hace este recorrido:

\begin{center}
Turret $\rightarrow$ Killjoy $\rightarrow$ Fade $\rightarrow$ Viper $\rightarrow$ Omen $\rightarrow$ Sage $\rightarrow$ Brimstone $\rightarrow$ Location server
\end{center}

Pasa por el AS 10, el AS 20, el AS 30, el AS 40 y llega al AS 50. En la traza los saltos aparecen como 172.16.120.1, 6.0.0.2, 9.0.0.1, 10.0.0.1, 11.0.0.2, 13.0.0.2 y 192.168.50.2. Cada router responde el ICMP \textit{Time Exceeded} con la IP de la interfaz por la que envía esa respuesta de vuelta a Turret. Por eso Viper aparece como 9.0.0.1 (su enlace hacia Omen) y no como 7.0.0.2, que es la interfaz por donde le entró el paquete. Lo mismo pasa con Omen (10.0.0.1) y Sage (11.0.0.2). La vuelta no repite el camino de ida: Sage le contesta a Turret por Reyna, y desde ahí la respuesta sigue por Gekko, Neon y Raze. Más abajo se explica por qué.

\subsubsection*{Por qué Killjoy sale directo por Fade}

Killjoy es router de frontera del AS 10 y conoce 192.168.50.0/24 por dos fuentes. La primera es eBGP, porque Fade se la anuncia directamente. La segunda es OSPF, porque Raze también la aprende por su lado (desde Neon) y la mete al AS 10 como ruta externa. Cuando dos protocolos distintos ofrecen el mismo prefijo, el router no compara métricas, porque no se pueden comparar entre protocolos. Elige por distancia administrativa: eBGP tiene 20 y OSPF tiene 110. Gana eBGP, y en la tabla de Killjoy la ruta aparece como \texttt{B 192.168.50.0/24 [20/0] via 6.0.0.2}.

\evid{Q3_ospf_vs_bgp/show_ip_route_killjoy.png}{Tabla de rutas de Killjoy. La ruta a Location server se aprende por eBGP desde Fade.}

Por eso el paquete de Turret ni siquiera recorre el interior del AS 10. Llega a su gateway, que ya es el router de salida, y de ahí cruza directo al AS 20.

\subsubsection*{Qué decide OSPF dentro del AS 10}

Para los routers internos, Deadlock y Cypher, Location server es una ruta externa. Raze y Killjoy la inyectan con \cmd{redistribute bgp 10 subnets} como tipo E2, con la misma métrica externa (20). Una E2 no suma el costo del camino interno, así que las dos opciones llegan empatadas. En IOS real ese empate se rompe con el costo hasta cada router de salida (ASBR). Para Cypher, Killjoy está a 64 (un serial) y Raze a 65 (Gigabit más serial), así que debería salir por Killjoy. Para Deadlock pasa lo contrario: Raze está a 64 y Killjoy a 65.

En Packet Tracer las tablas de Deadlock y Cypher muestran los dos caminos con métrica 20 y \textit{traffic share count} 1 en ambos. Es decir, el simulador no aplica ese desempate y reparte el tráfico entre los dos ASBR. Es una simplificación de Packet Tracer frente al comportamiento de un router real.

\evid{Q3_ospf_vs_bgp/route_location_cypher.png}{Ruta de Cypher a Location server: dos caminos externos, por Raze y por Killjoy.}
\evid{Q3_ospf_vs_bgp/route_location_deadlock.png}{Ruta de Deadlock a Location server.}

\subsubsection*{Qué decide BGP fuera del AS 10}

BGP no mira anchos de banda ni costos de enlace. Entre varias rutas a un mismo prefijo, y con los demás atributos iguales, prefiere la de AS-Path más corto, es decir, la que cruza menos sistemas autónomos. Desde el AS 10 hay dos salidas hacia Location server y las dos atraviesan los mismos AS: por Raze el camino es Neon, Gekko y Reyna (AS 20, 30 y 40), y por Killjoy es Fade, Viper y Omen (otra vez AS 20, 30 y 40). Como el largo empata, el AS-Path no es lo que define la salida. La definen la distancia administrativa en los routers de frontera y OSPF en los internos.

Donde sí se nota el AS-Path es en la vuelta. Packet Tracer no soporta iBGP (``Only external neighbors are supported''), así que dentro del AS 40 usamos OSPF 40 y Omen, Reyna y Sage redistribuyen entre OSPF y BGP. Cada vez que una ruta pasa por OSPF y vuelve a BGP pierde su AS-Path original. Por eso Location server se anuncia hacia afuera con camino \texttt{40} en vez de \texttt{40 50}, y las redes del AS 10 llegan a Viper desde Omen con camino \texttt{40}, que es más corto que \texttt{20 10} por Fade. Viper prefiere entonces devolver el tráfico por Omen, y eso produce el camino de vuelta distinto que se ve en la traza.

\subsubsection*{Resumen de la comparación}

OSPF elige dentro de un mismo sistema autónomo y usa una métrica técnica, el costo, que sale del ancho de banda de cada interfaz. Busca el camino más rápido según esa cuenta. BGP elige entre sistemas autónomos y usa el AS-Path, que solo cuenta cuántos AS se cruzan, sin saber nada de la velocidad de los enlaces. Cuando ambos ofrecen el mismo destino, la distancia administrativa decide cuál se instala en la tabla. En esta prueba eso hizo que Killjoy usara su vecino eBGP en vez de mandar el tráfico por dentro del AS 10.

Durante la integración aparecieron dos comportamientos de Packet Tracer que vale la pena dejar anotados. El primero es que, con \cmd{redistribute ospf 1} a secas, Raze y Killjoy reexportaban a BGP las rutas externas que habían aprendido de Omega. Neon y Fade las preferían porque tenían AS-Path \texttt{10}, y eso produjo un bucle entre Neon y Raze. Se corrigió con \cmd{redistribute ospf 1 match internal}, que es lo que IOS real hace por defecto. El segundo es que Killjoy instala como rutas B sus propias redes del AS 10, por ejemplo \texttt{172.16.24.0/21 [20/66] via 3.0.0.1}, después de redistribuirlas en BGP. El siguiente salto es el mismo que daría OSPF, así que el reenvío no cambia, pero en un router real esas redes aparecerían como O.

\subsection{Pregunta 4: tolerancia a fallos y reconvergencia}
\autor{Benjamín}

\textbf{Escenario:} comunicación entre un host de la LAN de Raze (Paint Shells, \texttt{172.16.64.2}) y uno de la LAN de Reyna (Dismiss, \texttt{183.112.0.2}). Se suspende la interfaz del enlace 5.0.0.0 (Raze--Neon) y se comprueba el camino alternativo por el enlace 6.0.0.0 (Killjoy--Fade). Al terminar se restaura el enlace. \pend{confirmar las direcciones de los hosts con \cmd{ipconfig} si el criterio de subredes cambia}

\subsubsection*{Ruta original}

Desde Paint Shells se ejecutó \cmd{tracert 183.112.0.2}, y en Raze se consultó la ruta hacia la LAN de Reyna.

\evid{Q4_caida_enlace/ruta_original.png}{Traza de Paint Shells a Dismiss con todos los enlaces activos.}
\evid{Q4_caida_enlace/show_ip_route_raze_antes.png}{Tabla de rutas de Raze hacia la LAN de Reyna antes de la caída.}

\pend{ruta activa router por router, y el protocolo (B u O) con que Raze aprendió la ruta}

\subsubsection*{Caída del enlace 5.0.0.0}

Se lanzó un ping continuo desde Paint Shells (\cmd{ping -t 183.112.0.2}) y se suspendió la interfaz del enlace en Raze:

\begin{lstlisting}
Raze# configure terminal
Raze(config)# interface GigabitEthernet0/1
Raze(config-if)# shutdown
\end{lstlisting}

\evid{Q4_caida_enlace/ping_continuo_caida.png}{Ping continuo durante la caída: paquetes perdidos hasta que se restablece la comunicación.}

\subsubsection*{Nueva ruta y tiempo de convergencia}

\evid{Q4_caida_enlace/ruta_alternativa.png}{Traza tras suspender el enlace 5.0.0.0.}
\evid{Q4_caida_enlace/show_ip_route_raze_despues.png}{Tabla de rutas de Raze tras la caída.}

\pend{nueva ruta, número de paquetes perdidos y tiempo de convergencia observado (cada ping sin respuesta equivale a cerca de un segundo), y justificación: detección inmediata de la caída de interfaz, recálculo en OSPF dentro del AS 10 y selección de otra ruta en BGP/OSPF}

\subsubsection*{Restauración}

\begin{lstlisting}
Raze(config-if)# no shutdown
\end{lstlisting}

\pend{confirmar con ping y traza que el tráfico volvió a la ruta original, y adjuntar la captura}

\subsection{Pregunta 5: impacto del ancho de banda en las métricas OSPF}
\autor{Sebastián}

\textbf{Escenario:} en el enlace Deadlock--Raze (2.0.0.0) se configura \cmd{bandwidth 150} y se envía un paquete de Annihilation a Paint Shells.

El costo OSPF de una interfaz es
\[
  \text{Costo} = \frac{10^{8}}{\text{Bandwidth (bps)}}
\]
con parte entera y ancho de banda de referencia de 100~Mbps, el valor por defecto de Cisco.

\subsubsection*{Cálculo de costos}
\begin{itemize}[nosep]
  \item Enlace directo con \cmd{bandwidth 150}: $\left\lfloor \dfrac{10^{8}}{150\,000} \right\rfloor = 666$, más el costo de la interfaz LAN de Raze (\pend{valor}).
  \item Trayecto alternativo Deadlock $\rightarrow$ Cypher $\rightarrow$ Killjoy $\rightarrow$ Raze: \pend{suma de los costos de cada enlace, según \cmd{show ip ospf interface}}.
\end{itemize}

\evid{Q5_ancho_banda/show_ip_route_deadlock.png}{Tabla de rutas de Deadlock con el ancho de banda modificado.}
\evid{Q5_ancho_banda/tracert_annihilation_paint.png}{Traza de Annihilation a Paint Shells.}

\subsubsection*{Conclusión}
\pend{¿el tráfico sigue por 2.0.0.0 o conmuta al trayecto alternativo? Corroborar con la tabla de rutas}

\textit{Recordatorio: restaurar el ancho de banda original después de la prueba.}

\subsection{Pregunta 6: auditoría de tablas de enrutamiento y redistribución}
\autor{Bárbara (Sage y distancias administrativas); Sebastián aporta la salida de Killjoy}

\subsubsection*{Killjoy}
\evid{Q6_tablas_rutas/show_ip_route_killjoy.png}{Salida de \cmd{show ip route} en Killjoy.}
\pend{señalar rutas O, O E2/E1 y B}

\subsubsection*{Sage}
\evid{Q6_tablas_rutas/show_ip_route_sage.png}{Salida de \cmd{show ip route} en Sage.}
\pend{señalar rutas B y las demás presentes}

\subsubsection*{Distancia administrativa}

La distancia administrativa (AD) es la confiabilidad que el router asigna a la fuente de una ruta; cuando dos protocolos aprenden el mismo prefijo, instala el de menor AD. Valores por defecto de Cisco IOS:

\begin{center}
\begin{tabular}{@{}lc@{}}
  \toprule
  Origen de la ruta & AD \\
  \midrule
  Red conectada (C) & 0 \\
  Ruta estática (S) & 1 \\
  eBGP (B) & 20 \\
  OSPF (O, O E1, O E2) & 110 \\
  iBGP (B) & 200 \\
  \bottomrule
\end{tabular}
\end{center}

\pend{comentar cuáles de estos valores aparecen entre corchetes en las capturas y qué efecto tienen en Killjoy, donde hay rutas OSPF y BGP}


\section{Conclusiones técnicas}

Análisis crítico del equipo sobre el comportamiento de OSPF y BGP, las dificultades encontradas y las lecciones aprendidas.

\subsection*{Sebastián}
\pend{párrafo sobre OSPF, redistribución y dificultades}

\subsection*{Benjamín}
\pend{párrafo sobre DHCP, BGP y reconvergencia}

\subsection*{Bárbara}
\pend{párrafo sobre BGP en el AS 40, tablas de rutas y distancia administrativa}

\subsection*{Visión conjunta}
\pend{síntesis del grupo: qué comportamiento de los protocolos sorprendió más y qué se haría distinto}


\end{document}
